%!TEX root = ../csuthesis_main.tex
\keywordsen{binocular vision;~3D human pose estimation;~lightweight network;~mixed-precision quantization;~mobile deployment}
\categoryen{TP391}
%\itemcountcn{There are \totalfigures\ figures, \totaltables\ tables, and \total{citnum}\ citations in this thesis.}
%\addcontentsline{toc}{section}{ABSTRACT}
\begin{abstracten}\setlength{\baselineskip}{20pt}
Multi-person 3D human pose estimation for mobile robotics, intelligent surveillance, and immersive interaction requires both spatial localization accuracy and real-time on-device inference. Monocular methods are limited by scale and depth ambiguity, while existing stereo methods often suffer from redundant dual-branch feature extraction, dependence on dense disparity or cost-volume computation, and poor efficiency on mobile hardware. To address these issues, the main work of this thesis can be summarized in the following two aspects:

  (1) To address redundant computation and error accumulation in on-device stereo methods, this thesis proposes a lightweight binocular 3D human pose estimation framework for multi-person on-device scenarios. The framework jointly redesigns the stereo feature extraction, cross-view interaction, and explicit geometric solving pipeline, replacing the traditional cascade paradigm based on dense disparity estimation with an end-to-end lightweight design. On the RT-Pose benchmark, the proposed method achieves 158.6 mm MPJPE, reducing error by 48.4\% over the best monocular baseline and by 21.5\% over the best cascade baseline, while consuming only 13.93 G FLOPs---88.0\% fewer than the cascade baseline relying on learned stereo depth. In the zero-shot generalization test on MADS, the model achieves 33.68 mm MPJPE, demonstrating robustness to cross-domain scenarios.

  (2) To address latency and quantization sensitivity during on-device deployment, this thesis further proposes a hardware-aware mixed-precision deployment scheme. The scheme applies end-to-end latency breakdown, Roofline analysis, and module-level sensitivity screening to identify backbone convolutions as the primary quantization target and cross-view interaction and geometric solving modules as high-precision protection paths, yielding a layer-wise mixed-precision quantization plan. On Xiaomi\,14, the scheme achieves a 1.60$\times$ inference speedup and reduces model size by 35.0\%, with only a 1.4\% increase in MPJPE. On the Jetson AGX Xavier + TensorRT platform, the mixed-precision engine reaches 48.5 FPS while reducing GPU memory by 27.7\% and per-frame energy consumption by 14.9\%, again with no more than 1.4\% MPJPE degradation, demonstrating the portability and practicality of the proposed scheme on embedded GPU hardware.

\end{abstracten}
